La filosofía del DCU, presentada en el capítulo anterior, establece que las decisiones de diseño deben fundamentarse en las necesidades reales de quienes interactúan con el producto. Este principio se adoptó en el rediseño de los manipuladores de Justina aplicando la metodología de Karl T.\ Ulrich y Steven D.\ Eppinger.

Ulrich y Eppinger presentan en \textit{Diseño y Desarrollo de Productos} un conjunto de métodos para el diseño de productos físicos, discretos e ingenieriles~\cite{Ulrich-Eppinger}. Su proceso genérico se divide en seis fases: planeación, desarrollo del concepto, diseño a nivel sistema, diseño de detalle, pruebas y refinamiento, e inicio de producción.

Este capítulo abarca las fases de desarrollo del concepto y diseño a nivel sistema, que comprenden la identificación de necesidades, el establecimiento de especificaciones, la descomposición funcional, la generación de conceptos y la selección del concepto que se llevará al diseño de detalle.

\section{Identificación de necesidades}

La identificación de necesidades del cliente es la fase fundacional del proceso. Su objetivo es capturar de manera sistemática y exhaustiva las necesidades, expectativas y restricciones de los usuarios potenciales, transformando información cualitativa en un conjunto estructurado de requerimientos. Ulrich y Eppinger~\cite{Ulrich-Eppinger} definen una necesidad del cliente como una descripción, en las propias palabras del usuario, de los beneficios que un producto debe proporcionar, y distinguen con claridad entre necesidades ---que describen \emph{qué} debe lograr el producto--- y soluciones ---que especifican \emph{cómo} lograrlo---. El proceso se estructura en cinco pasos: recopilación de datos sin procesar mediante entrevistas, interpretación de esos datos en términos de necesidades, organización jerárquica, establecimiento de importancia relativa y reflexión sobre los resultados.

\subsection{Perfiles de usuario}

Para la implementación del proceso de identificación de necesidades se definieron dos perfiles de participantes mediante muestreo intencional~\cite{Protocolo}.

El primero es el \textbf{usuario final}: personas que requieren asistencia doméstica constante debido a edad avanzada o alguna condición que limite su autonomía en actividades cotidianas del hogar, con capacidad cognitiva para participar en una entrevista y proporcionar consentimiento informado~\cite{Protocolo}. Sus necesidades se relacionan con la utilidad práctica del robot, la facilidad de interacción, la confiabilidad y la seguridad durante el uso.

El segundo perfil corresponde a \textbf{desarrolladores de software con experiencia en robótica de servicio}: miembros activos o previos del equipo de desarrollo de Justina con experiencia de participación en RoboCup@HOME y conocimiento técnico sobre el diseño, programación o implementación de funcionalidades del robot~\cite{Protocolo}. Este perfil aporta una perspectiva complementaria: han trabajado directamente con robots de servicio en condiciones desafiantes, han observado y comparado múltiples plataformas en competencia, y constituyen también un tipo de usuario del sistema al ser quienes realizan el mantenimiento y la programación del robot.

Dentro de este segundo perfil se distinguió entre quienes trabajan con \textbf{plataformas experimentales} de diseño propio y quienes utilizan \textbf{plataformas comerciales}, específicamente el \textit{Toyota Human Support Robot} (Takeshi) del Laboratorio de Bio-Robótica~\cite{Protocolo}. Esta distinción motivó el diseño de variantes específicas en los instrumentos de recolección.

\subsection{Recopilación de datos}

El proyecto contó con la aprobación del Comité de Ética en Investigación y Docencia de la Facultad de Ingeniería de la UNAM~\cite{Consentimiento}. Todos los participantes firmaron una carta de consentimiento informado antes de su participación, bajo la cual se garantizó la confidencialidad de la información mediante codificación de datos, almacenamiento seguro y acceso restringido al equipo investigador~\cite{Consentimiento}.

Siguiendo las recomendaciones de Ulrich y Eppinger~\cite{Ulrich-Eppinger}, se seleccionaron las entrevistas como método principal de recopilación. Se diseñaron cuatro guiones semiestructurados diferenciados según perfil: uno para el usuario final, uno de aplicación general para desarrolladores, y dos variantes específicas para desarrolladores de plataformas experimentales y comerciales, respectivamente~\cite{Protocolo}. Los guiones completos se documentan en el Anexo~[X].

Se realizaron nueve entrevistas en total: cinco con desarrolladores con experiencia en RoboCup@HOME y cuatro con usuarios finales potenciales. Todas fueron grabadas en audio con consentimiento explícito de los participantes~\cite{Consentimiento}. Los fragmentos relevantes se transcribieron y organizaron en tablas estructuradas que vinculan cada declaración con su participante, el contexto de la pregunta y anotaciones sobre necesidades implícitas.

\subsection{Interpretación y organización jerárquica}

Siguiendo las directrices de Ulrich y Eppinger~\cite{Ulrich-Eppinger}, cada necesidad se expresó como un atributo del producto ---sujeto, verbo que describe la acción o capacidad deseada y, cuando es necesario, objeto o contexto de aplicación--- capturando el \textit{qué} sin anticipar el \textit{cómo}. Se adoptó un enfoque holístico que abarcó el robot de servicio doméstico en su totalidad, reconociendo que las necesidades de manipulación no existen de manera aislada sino interrelacionadas con navegación, percepción e interacción. Como resultado, se identificaron \textbf{122 necesidades distintas}, mencionadas colectivamente \textbf{690 veces} a lo largo de todas las entrevistas.

Las 122 necesidades se organizaron en una jerarquía de \textbf{14 necesidades primarias} con sus correspondientes necesidades secundarias y terciarias. El agrupamiento se realizó en una hoja de cálculo de Excel, lo que permitió reorganizar iterativamente los enunciados y mantener trazabilidad entre cada necesidad y su fuente original. La estructura jerárquica completa se documenta en el Anexo~[Y].

\subsection{Importancia relativa y selección}

Para determinar la importancia relativa de las 14 necesidades primarias se implementó una estrategia combinada. El primer componente fue la \textbf{frecuencia relativa de mención}: el cociente entre el número de menciones de cada necesidad primaria y sus derivadas sobre el total de 690 menciones. El segundo componente fue una \textbf{encuesta estructurada} distribuida electrónicamente a tres de los participantes entrevistados, cuyos resultados se interpretaron como indicadores cualitativos de tendencia complementarios a los datos de frecuencia, dado el tamaño reducido de la muestra. La escala utilizada constó de cinco niveles: \textit{Imprescindible} (1.0), \textit{Deseable} (0.6), \textit{Indiferente} (0.4), \textit{Preferiría que no la tuviera} (0.2) y \textit{No debería tenerla} (0.0). La importancia relativa final se calculó combinando el promedio de las evaluaciones directas con la frecuencia de mención de cada necesidad.

Aplicando un umbral de corte basado en la distribución de valores obtenidos, se seleccionaron seis necesidades primarias prioritarias. De estas, se retuvieron tres con implicación directa en el diseño de los manipuladores: \textit{el robot manipula artículos domésticos}, \textit{el robot es robusto} y \textit{el robot tiene un diseño intuitivo}. A la primera se le incorporaron cuatro necesidades secundarias que detallan aspectos específicos de la capacidad de manipulación. El conjunto final quedó conformado por:

\begin{enumerate}
    \item \textbf{El robot manipula artículos domésticos.}
    \begin{itemize}
        \item El robot planea trayectorias del manipulador.
        \item El robot toma artículos con la fuerza necesaria.
        \item El robot extiende el alcance efectivo de sus manipuladores.
        \item El robot manipula equipo para realizar tareas domésticas.
    \end{itemize}
    \item \textbf{El robot es robusto.}
    \item \textbf{El robot tiene un diseño intuitivo.}
\end{enumerate}

% ============================================================
%  SECCIÓN 2: ESPECIFICACIONES DEL PRODUCTO
% ============================================================

\section{Establecimiento de especificaciones objetivo}

La segunda fase de la metodología consiste en traducir las necesidades del cliente en especificaciones técnicas precisas y medibles. Ulrich y Eppinger~\cite{Ulrich-Eppinger} distinguen entre necesidades ---expresadas en el lenguaje subjetivo del usuario--- y especificaciones ---que describen con exactitud cuantificable qué debe ser el producto---. Una especificación consiste en una \textit{métrica} y un \textit{valor} acompañado de sus unidades. Las especificaciones se establecen en dos momentos: primero como \textbf{especificaciones objetivo}, antes de conocer las restricciones técnicas del concepto de diseño, y luego como \textbf{especificaciones finales}, una vez seleccionado el concepto. El proceso comprende cuatro pasos: elaborar la lista de métricas, recabar información de referencia mediante \textit{benchmarking}, establecer valores meta ideales y marginalmente aceptables, y reflexionar sobre los resultados.

\subsection{Traducción de necesidades a métricas}

Para cada necesidad del cliente se identificó la característica precisa y medible del producto que mejor refleja el grado en que esa necesidad queda satisfecha. La relación entre necesidades y métricas se documentó en una matriz de correlación (Cuadro~\ref{tab:correlacion}) que evidencia qué métricas cubren cada necesidad y permite identificar aquellas que requieren múltiples indicadores. A partir de las necesidades identificadas se definieron las ocho métricas presentadas en el Cuadro~\ref{tab:metricas}.

\begin{table}[htbp]
\centering
\caption{Lista de métricas derivadas de las necesidades del cliente.}
\label{tab:metricas}
\begin{tabular}{clccc}
\hline
\textbf{\#} & \textbf{Métrica} & \textbf{Nec. relacionadas} & \textbf{Importancia} & \textbf{Unidad} \\
\hline
1 & Capacidad de carga     & 1, 5, 6 & 5 & {[}kg{]} \\
2 & Fuerza de agarre       & 1, 3, 5 & 5 & {[}N{]}  \\
3 & Alcance                & 2, 4, 5 & 5 & {[}mm{]} \\
4 & Apertura               & 1, 5    & 4 & {[}mm{]} \\
5 & Grados de libertad     & 4, 5    & 4 & {[}1{]}  \\
6 & Repetibilidad          & 2, 5, 6 & 4 & {[}mm{]} \\
7 & Grado de protección IP & 5, 6    & 3 & {[}1{]}  \\
8 & Mapa de desensamble    & 6, 7    & 3 & {[}s{]}  \\
\hline
\end{tabular}
\end{table}

\begin{table}[htbp]
\centering
\caption{Matriz de correlación entre necesidades del cliente y métricas.}
\label{tab:correlacion}
\small
\begin{tabular}{lcccccccc}
\hline
\textbf{Necesidad} & \textbf{M1} & \textbf{M2} & \textbf{M3} & \textbf{M4} & \textbf{M5} & \textbf{M6} & \textbf{M7} & \textbf{M8} \\
\hline
El robot manipula artículos domésticos        & $\times$ & $\times$ &          & $\times$ &          &          &          &          \\
El robot planea trayectorias del manipulador  &          &          & $\times$ &          &          & $\times$ &          &          \\
El robot toma artículos con la fuerza nec.    &          & $\times$ &          &          &          &          &          &          \\
El robot extiende el alcance efectivo         &          &          & $\times$ &          & $\times$ &          &          &          \\
El robot manipula equipo para tareas dom.     & $\times$ & $\times$ & $\times$ & $\times$ & $\times$ & $\times$ & $\times$ &          \\
El robot es robusto                           & $\times$ &          &          &          &          & $\times$ & $\times$ & $\times$ \\
El robot tiene un diseño intuitivo            &          &          &          &          &          &          &          & $\times$ \\
\hline
\end{tabular}
\end{table}

\subsection{Benchmarking y análisis competitivo}

Para establecer valores de referencia se realizó un análisis comparativo de sistemas representativos de dos categorías: \textit{cobots} de uso industrial ---Panda/Franka~Emika, Gen3/Kinova, UR3e/Universal~Robots, EBLM/Ti5~Robot y OpenArm--- y manipuladores de robots de servicio ---Stretch~3/Hello~Robot, HSR/Toyota, TIAGo/PAL~Robotics y G1/Unitree---. Para los efectores finales se compararon grippers de dos dedos (Franka~Hand, 2F-85 y 2F-140/Robotiq, HSR) y de tres dedos (3F-155/Robotiq, 3F-AR/UW-Milwaukee, DG-3F-B/Tesollo, Dex3s/Unitree).

La evaluación se realizó mediante un análisis de calidad de función (QFD) que ponderó cada métrica según su importancia relativa a las necesidades primarias identificadas. Para el brazo se evaluaron cuatro grupos de necesidades: planeación de trayectorias (alcance y repetibilidad), alcance efectivo (alcance y grados de libertad), tareas domésticas (capacidad de carga, alcance, grados de libertad, repetibilidad e IP) y robustez (capacidad de carga, repetibilidad e IP). Para el gripper se evaluaron las necesidades de manipulación (capacidad de carga, fuerza de agarre, apertura, grados de libertad y repetibilidad), control de fuerza, tareas domésticas y robustez.

\subsection{Definición de valores objetivo}

Con base en el análisis de referencia se establecieron valores objetivo para cada métrica en dos niveles: el \textbf{valor ideal}, que representa el mejor desempeño razonablemente alcanzable, y el \textbf{valor marginal}, que constituye el mínimo aceptable para que el diseño sea viable. El Cuadro~\ref{tab:especificaciones} presenta el conjunto completo de especificaciones objetivo.

\begin{table}[htbp]
\centering
\caption{Especificaciones objetivo para los manipuladores de Justina.}
\label{tab:especificaciones}
\begin{tabular}{clcccc}
\hline
\textbf{\#} & \textbf{Métrica} & \textbf{Imp.} & \textbf{Unidad} & \textbf{V. marginal} & \textbf{V. ideal} \\
\hline
1 & Capacidad de carga     & 5 & {[}kg{]}  & $>2$    & $>3$      \\
2 & Fuerza de agarre       & 5 & {[}N{]}   & $>70$   & $>100$    \\
3 & Alcance                & 5 & {[}mm{]}  & $>500$  & $>700$    \\
4 & Apertura               & 4 & {[}mm{]}  & $>100$  & $>140$    \\
5 & Grados de libertad     & 4 & {[}1{]}   & $\geq5$ & $\geq6$   \\
6 & Repetibilidad          & 4 & {[}mm{]}  & $<\pm1$ & $<\pm0.1$ \\
7 & Grado de protección IP & 3 & {[}1{]}   & IP44    & IP55      \\
8 & Mapa de desensamble    & 3 & {[}s{]}   & ---     & ---       \\
\hline
\end{tabular}
\end{table}

La métrica M8 (mapa de desensamble), correspondiente a la necesidad de diseño intuitivo, se retiene en el conjunto como requerimiento cualitativo: su naturaleza no admite una cuantificación directa en esta etapa y su evaluación se formalizará durante la fase de selección de conceptos.

% ============================================================
%  SECCIÓN 3: GENERACIÓN DE CONCEPTOS
% ============================================================

\section{Generación de conceptos}

La generación de conceptos es la actividad mediante la cual el equipo de desarrollo explora de manera sistemática el espacio de alternativas de diseño posibles antes de comprometerse con una solución particular. Ulrich y Eppinger~\cite{Ulrich-Eppinger} definen el concepto de un producto como una descripción aproximada de la tecnología, los principios de funcionamiento y la forma del producto. El método estructurado propuesto consta de cinco pasos: aclarar el problema mediante descomposición funcional, buscar externamente conceptos existentes, generar internamente nuevas soluciones, explorar sistemáticamente las combinaciones posibles y reflexionar sobre los resultados. Dado que la generación de conceptos es relativamente económica en comparación con las fases posteriores del desarrollo, los autores argumentan que no existe justificación para limitarla, y que un equipo eficiente debería explorar el espacio de alternativas con amplitud suficiente como para tener confianza en que no ha pasado por alto soluciones superiores.

\subsection{Descomposición funcional del problema}

El primer paso consistió en representar el sistema de manipulación como una caja negra que opera sobre flujos de energía, material y señales, para luego dividirla en subfunciones que describen lo que los elementos del manipulador deben hacer. De la totalidad de subfunciones identificadas se seleccionaron los \textbf{seis subproblemas críticos} que determinan en mayor medida el desempeño del sistema: aceptar energía, convertir energía a movimiento, transmitir movimiento y fuerza, sensar posición, sensar fuerza y sostener el objeto. El diagrama funcional completo, que muestra la interacción entre estas subfunciones a través de los flujos de energía y señal, se presenta en la Figura%~\ref{fig:diagrama_funcional}.

%\begin{figure}[htbp]
%    \centering
%    \includegraphics[width=\textwidth]{diagrama_funcional}
%    \caption{Diagrama funcional del sistema de manipulación de Justina. Las líneas continuas representan flujos de material y energía; las líneas discontinuas representan señales de control y retroalimentación.}
%    \label{fig:diagrama_funcional}
%\end{figure}

\subsection{Búsqueda externa de soluciones}

Para cada uno de los seis subproblemas críticos se realizó una búsqueda exhaustiva de soluciones existentes en la literatura técnica, catálogos de fabricantes, patentes y en el análisis directo de los robots estudiados durante el benchmarking. Las soluciones encontradas se organizaron en tablas morfológicas ---una por subproblema--- que cubren los principales principios de solución para cada función. Las tablas completas se incluyen en el Anexo~[W].

\subsection{Exploración sistemática y acotamiento del espacio de soluciones}

La combinación irrestricta de las soluciones identificadas genera un espacio de diseño de gran magnitud. Mediante un script de Python se enumeraron sistemáticamente \textbf{1943 combinaciones} que representan el espacio de búsqueda completo antes de aplicar criterios de compatibilidad.

Para reducir este espacio a candidatos viables se procedió en dos etapas. Primero, se acotó el conjunto de opciones por subproblema a las soluciones más pertinentes para el contexto de aplicación: robot de servicio doméstico, operación eléctrica e integración con el stack de software existente del laboratorio. Segundo, sobre este espacio reducido se aplicó un criterio de \textbf{afinidad entre soluciones}: se conservaron únicamente aquellas combinaciones que representan configuraciones lógicamente coherentes y directamente implementables, que no requieren adaptaciones innecesarias para funcionar conjuntamente. Este filtrado redujo el espacio de candidatos a \textbf{384 combinaciones realistas}.

\subsection{Desarrollo de conceptos completos}

Del conjunto de 384 candidatos se seleccionaron \textbf{4 conceptos representativos} para su desarrollo en detalle, buscando cubrir la diversidad de enfoques presentes en el espacio de soluciones. Cada concepto se materializó en un modelo CAD conceptual que permitió evaluar de manera preliminar las métricas establecidas en las especificaciones objetivo y facilitar la comparación estructurada en la fase de selección. Los cuatro conceptos, sus características principales y su evaluación frente a las especificaciones se describen en el capítulo siguiente.
